\subsection{A finite-checkpoint confidence certificate for independent shift batches}
\label{sec:rqmc-confidence-certificate}

I treat an entire independently shifted digital net as one observation. I assume
that $B_1,\ldots,B_R$ are independent copies of a variable $B\in[0,\rho]$, where
$0\leq\rho\leq1$. I write $\mu=\mathbb E B$, $\sigma^2=\operatorname{Var}(B)$,
$m=\lfloor R/2\rfloor$, and
\[
 Z_i=\tfrac12(B_{2i-1}-B_{2i})^2,\quad
 M=\rho^2/2,\quad
 \overline Z=m^{-1}\sum_{i=1}^m Z_i,\quad
 \overline B=R^{-1}\sum_{j=1}^R B_j.
\]
The $Z_i$ are independent, $0\leq Z_i\leq M$, and
$\mathbb E Z_i=\sigma^2$. The last identity follows by expanding the square
and using independence of the two batch means in each pair. It does not require
independence among financial evaluations within a batch.

\begin{lemma}[A pair-based upper confidence bound on the batch variance]
\label{lem:rqmc-pair-variance}
Let $0<\delta<1$, put $t=\log(1/\delta)/m$, and define
\[
 U_{\rm pair}=\overline Z+Mt+
       \sqrt{M^2t^2+2Mt\overline Z}.
\]
Then $\mathbb P\{\sigma^2>U_{\rm pair}\}\leq\delta$.
If a deterministic certificate gives $\sigma^2\leq V_*$, the same statement
holds for $U=\min\{V_*,U_{\rm pair}\}$.
\end{lemma}

\begin{proof}
For $x\geq0$, $e^{-x}\leq1-x+x^2/2$. One proof observes that
$g(x)=1-x+x^2/2-e^{-x}$ has $g(0)=g'(0)=0$ and
$g''(x)=1-e^{-x}\geq0$. Since $Z_i^2\leq MZ_i$, for every $\theta\geq0$,
\[
 \mathbb E e^{-\theta Z_i}
 \leq1-\theta\sigma^2+\theta^2M\sigma^2/2
 \leq\exp(-\theta\sigma^2+\theta^2M\sigma^2/2).
\]
Multiplication over the independent pairs and Markov's inequality give, for
$s\geq0$ and $\sigma^2>0$,
\[
 \mathbb P\{\overline Z\leq\sigma^2-s\}
 \leq\inf_{\theta\geq0}
       \exp(-m\theta s+m\theta^2M\sigma^2/2)
 =\exp\!\left(-\frac{ms^2}{2M\sigma^2}\right).
\]
When $M=0$ or $\sigma^2=0$, all $Z_i$ vanish almost surely and the conclusion
is immediate. In the other cases $U_{\rm pair}$ is the larger root, as a
function of $v$, of $(v-\overline Z)^2=2Mtv$. Therefore
$\sigma^2>U_{\rm pair}$ implies
$\overline Z<\sigma^2-\sqrt{2Mt\sigma^2}$. If this threshold is negative,
the event is empty. Otherwise the preceding inequality, with
$s=\sqrt{2Mt\sigma^2}$, bounds its probability by $e^{-mt}=\delta$.
The deterministic prior can be imposed by intersection without introducing
another random event.
\end{proof}

\begin{proposition}[A mean interval with a variance estimate from the same observations]
\label{prop:rqmc-mean-confidence}
Let $0<\alpha_v,\alpha_m<1$, use the preceding lemma with
$\delta=\alpha_v$, and set
\[
 L=\log(2/\alpha_m),\qquad b=\frac{\rho L}{3R},\qquad
 w(U)=b+\sqrt{b^2+\frac{2UL}{R}}.
\]
Suppose that a deterministic mean certificate gives $0\leq\mu\leq u\leq\rho$.
The possibly empty interval
\[
 I=[\max\{0,\overline B-w(U)\},\ \min\{u,\overline B+w(U)\}]
\]
contains $\mu$ with probability at least $1-\alpha_v-\alpha_m$. An empty
intersection is an explicit confidence-failure outcome and cannot supply a
price estimate. The variance and mean observations need not be independent.
\end{proposition}

\begin{proof}
For completeness, put $Y=B-\mu$. Then $\mathbb E Y=0$, $|Y|\leq\rho$ and
$\mathbb E Y^2=\sigma^2$. For $k\geq2$,
$\mathbb E|Y|^k\leq\rho^{k-2}\sigma^2$ and
$k!\geq2\,3^{k-2}$. Expanding the exponential and bounding the absolute
values of its terms show that, for $|\theta|\rho<3$,
\[
 \mathbb E e^{\theta Y}
 \leq 1+\frac{\sigma^2\theta^2}{2(1-|\theta|\rho/3)}
 \leq\exp\!\left(\frac{\sigma^2\theta^2}{2(1-|\theta|\rho/3)}\right).
\]
This derivation applies to $Y$ and $-Y$. Independence, Markov's inequality
and the choice $\theta=e/(\sigma^2+\rho e/3)$ give, for $e>0$ and
$\sigma^2>0$,
\[
 \mathbb P\{|\overline B-\mu|\geq e\}
 \leq2\exp\!\left(-\frac{Re^2}{2\sigma^2+2\rho e/3}\right).
\]
If $\sigma^2=0$, $\overline B=\mu$ almost surely and the required conclusion
is immediate. Solving $Re^2=L(2\sigma^2+2\rho e/3)$ gives $e=w(\sigma^2)$.
Thus the event
$E_m=\{|\overline B-\mu|\leq w(\sigma^2)\}$ has probability at least
$1-\alpha_m$. The pair lemma gives
$E_v=\{\sigma^2\leq U\}$ with probability at least $1-\alpha_v$.
On their intersection, monotonicity of $w$ gives
$|\overline B-\mu|\leq w(U)$, and the deterministic prior contains $\mu$.
A union bound therefore proves the stated coverage. In particular, replacing
$\sigma^2$ by the random $U$ is justified on a joint event, rather than by
conditioning a Bernstein inequality on an estimated variance.
\end{proof}

I prescribe eleven checkpoints $R_j=2^{6+j}$, $j=0,\ldots,10$, before
observing any shifts. At every checkpoint I assign
$\alpha_v=\alpha_m=1/2200$. Applying the proposition at each checkpoint
and then a union bound gives simultaneous coverage at least
$1-22/2200=0.99$ for one selected price. This permits stopping at the first
of these checkpoints whose nonempty interval meets the retained tolerance.
It does not establish coverage for an unprescribed stopping time, or
99 per cent simultaneous coverage for all 27 prices. The batches at different
checkpoints share observations; independence between checkpoints is unnecessary
for this finite union bound.

The numerical implementation replaces both logarithms by upward enclosures.
For rational arguments $1<x\leq10^{40}$, I evaluate their quotient and logarithm
with correctly rounded \texttt{Decimal} arithmetic at 160 significant digits,
and add $10^{-120}$. The quotient has relative rounding error at most
$5\times10^{-160}$; $\log x<93$, and the elementary logarithm and final addition
have absolute rounding errors below $10^{-157}$. The increment therefore
strictly exceeds the accumulated rounding error. A rational square-root
argument $a\geq0$ is enclosed by
\[
 \operatorname{sqrt}_{+}(a)=
  2^{-160}\left\lceil\sqrt{\lceil 2^{320}a\rceil}\right\rceil.
\]
Both integer roundings are evaluated with arbitrary-precision integers and
$\operatorname{sqrt}_{+}(a)^2\geq a$ exactly. These upward substitutions
can only enlarge $U_{\rm pair}$ and $w(U)$, apart from the justified cap by
$V_*$. Interval endpoints and recorded statistics are rational. The software
premise of correctly rounded elementary functions remains explicit.

The probability inequalities above are established statistical mechanisms,
with published empirical-Bernstein randomised quasi-Monte Carlo antecedents.
I use the derivation to justify the particular finite-checkpoint certificate,
not to claim a new generic confidence theorem. The separate retained audit
checks every stored integer statistic and endpoint, and exactly sums
Bernoulli pair-type probabilities at $R=64$ and $R=128$. These finite
checks provide implementation evidence; the proof supplies the distribution-free
claim under its stated independence and range hypotheses.
